\noindent\textbf{Elección del procedimiento.} Imponer $x^*=rx^*(1-x^*)$ da $x^*[r(1-x^*)-1]=0$. Hay un punto fijo nulo y otro $1-1/r$. Para perturbaciones pequeñas, $\delta x_{n+1}\simeq f\prime(x^*)\delta x_n$; estabilidad local requiere módulo del multiplicador menor que uno.

\begin{enumerate}

\item Dividir por el punto fijo no nulo para obtener la segunda raíz.

\[\fitmath{\mathrm{xs}=1-\frac{1}{r}}\]

\[\fitmath{\mathrm{xs}=1-\frac{1}{(2{,}8)}\approx 0{,}6428571\,}\]

\item Derivar $rx(1-x)$ y evaluar en $x^*=1-1/r$.

\[\fitmath{\mathrm{mult}=2-r}\]

\[\fitmath{\mathrm{mult}=2-(2{,}8)\approx -0{,}8\,}\]

\end{enumerate}

\noindent\textbf{Contraste y comparación.} El multiplicador -0,8 tiene módulo menor que uno: las perturbaciones pequeñas alternan de signo y decrecen. El punto fijo nulo tiene multiplicador 2,8 y es inestable.
